% !TITLE 敕勒歌
% !AUTHOR 无名氏
% !DYNASTY 南北朝
% !GENRE 北朝民歌
% !ORDER 2

\begin{poem}
敕勒川\textsuperscript{1}，阴山\textsuperscript{2}下。
天似穹庐\textsuperscript{3}，笼盖四野。
天苍苍\textsuperscript{4}，野茫茫\textsuperscript{5}，
风吹草低见\textsuperscript{6}牛羊。
\end{poem}

\begin{annotations}
\notetext{1}{敕勒川：敕勒族居住的草原。敕勒是南北朝时期北方的一个游牧民族，敕勒川大致在今天的蒙古草原一带。}
\notetext{2}{阴山：山脉名，横亘今天的内蒙古中部，自古以来是游牧民族与农耕民族的分界线。}
\notetext{3}{穹庐：游牧民族住的圆顶帐篷，即蒙古包。用帐篷比喻天空，是牧民独特的想象力——天就是最广阔的“家”。
}
\notetext{4}{苍苍：深青色。天色蓝得发青，一望无际，干净得什么也没有。}
\notetext{5}{茫茫：辽阔无边的样子。原野大得看不到尽头，人在其中渺小如一粒沙。}
\notetext{6}{见：同“现”，显现、显露。风吹过，草低伏下去，隐藏的牛羊就露了出来。这是一个动态的画面——你仿佛能看见草浪翻滚、牛羊出没的壮阔场景。}
\end{annotations}

\begin{appreciation}
全诗二十七字，是中国文学史上最短的写景诗之一，也是装下的天地最大的一首。短得连一条微博都填不满，却写出了一片完整的草原。

敕勒是南北朝时北方的游牧民族，敕勒川就是他们在阴山脚下放牧的草原。阴山横在今天内蒙古中部，自古是游牧和农耕的分界线——翻过这座山，日子就是另一种过法。这首歌原本用敕勒语唱，后来翻译成汉语，成了今天的样子。翻译得这么顺口，因为它本来就是民歌，本来就在人嘴里传来传去。

"天似穹庐，笼盖四野"——天像一顶大帐篷，笼罩着四面的原野。穹庐就是蒙古包。这个比喻，只有住在帐篷里的人想得出来：对牧民来说，天不是抽象的概念，而是一顶更大的帐篷，是家的延伸。我们读《诗经》时说过，好的比喻都来自真实的生活，这句就是活例子。

"天苍苍，野茫茫"：天是深蓝发青的，原野一望无边。记得《蒹葭》的"蒹葭苍苍"吗？同一个"苍苍"，那边写秋日芦苇的迷蒙，这边写天空的辽阔，各是各的味道。两个叠词念起来像呼吸，一呼一吸，天地之间就剩你一个人。

全诗的眼睛是最后一句："风吹草低见牛羊"。见通"现"，显露。草原上的草能没到马肚子，牛羊低头吃草，就藏进草浪里看不见。忽然一阵风，草伏下去，成群的牛羊露了出来。七个字，有风，有草，有牛羊，还是一个连贯的动态——风来，草伏，牛羊现。闭上眼，草浪就滚起来了。

北朝民歌和南朝民歌，像南北方的水土一样不同：南朝写"采莲南塘秋"，是水乡的精巧；北朝写草原骏马，是天地间的苍茫。可追到根上，它们都和《诗经》的国风一路：朴素，直接，不端着。民歌的底子都是这个。

妹妹，你还没亲眼见过草原吧？哪天放长假，我们去一趟，亲眼看看风吹草低见牛羊，比背一百遍都管用。
\end{appreciation}
